
\slajdid{05-s018}
\begin{frame}[label=05-s018]{Oduzimanje}
\gusttekst\setlength{\parskip}{1pt}
PRIMER:\quad $326.23_{10}-95.9_{10}={?}_{10}$
\begin{center}\def\DEoper{-}\def\DErezultat{2,3,0,.,3,3}% Parametri DEoper (+/-) i DErezultat. Svi prikazani zapisi ostaju zasebni.
\begin{tikzpicture}[x=.35cm,y=.45cm,font=\fontsize{9}{10}\selectfont]
\foreach \x/\v in {0/3,1/2,2/6,3/.,4/2,5/3}{\node[anchor=base] at (\x,1) {$\v$};}
\foreach \x/\v in {1/9,2/5,3/.,4/9}{\node[anchor=base] at (\x,0) {$\v$};}
\node[anchor=base] at (-1,0) {$\DEoper$};\draw[DEplava] (3.18,-.35)--(3.18,1.5);
\draw[DEplava,line width=2pt,-{Latex[length=3mm]}] (7.2,.5)--(10.2,.5);
\node[align=center,text width=4.3cm] at (2,-1.4) {USAGLAŠENO\\MESTO TAČAKA};
\begin{scope}[shift={(12,0)}]
\foreach \x/\v in {0/3,1/2,2/6,3/.,4/2,5/3}{\node[anchor=base] at (\x,1) {$\v$};}
\foreach \x/\v in {0/0,1/9,2/5,3/.,4/9,5/0}{\ifnum\x=0\def\DEboja{DEcrvena}\else\ifnum\x=5\def\DEboja{DEcrvena}\else\def\DEboja{DEtekst}\fi\fi\node[\DEboja,anchor=base] at (\x,0) {$\v$};}
\node[anchor=base] at (-1,0) {$\DEoper$};
\node[align=center,text width=5.1cm] at (2,-1.4) {PODRAZUMEVANE VODEĆE I\\PRATEĆE NULE};
\end{scope}
\begin{scope}[shift={(6,-4.3)}]
\foreach \x/\v in {0/3,1/2,2/6,3/.,4/2,5/3}{\node[anchor=base] at (\x,1) {$\v$};}
\foreach \x/\v in {0/0,1/9,2/5,3/.,4/9,5/0}{\ifnum\x=0\def\DEboja{DEcrvena}\else\ifnum\x=5\def\DEboja{DEcrvena}\else\def\DEboja{DEtekst}\fi\fi\node[\DEboja,anchor=base] at (\x,0) {$\v$};}
\node[anchor=base] at (-1,0) {$\DEoper$};\draw[DEplava] (-1.3,-.45)--(5.4,-.45);
\foreach \v [count=\x from 0] in \DErezultat {\node[anchor=base] at (\x,-1.15) {$\v$};}
\end{scope}
\end{tikzpicture}
\end{center}
Krećemo oduzimanje od cifara najmanje težine\\
1. $3-0=3$\\
2. $2-9={?}$ „Ne može“ pa ćemo pozajmiti jedinicu od cifara veće težine.\\
\hspace*{5mm}$12-9=3$ i pozajmica 1\\
3. $6-(5+1)$ ili $(6-1)-5$ kako smo već učili, pošto imamo pozajmicu jedan zbog oduzimanja cifara manje težine\\
\ldots\\
K. Tačku stavljamo na isto mesto koje je bilo u operandima
\end{frame}
